\section{Conclusion}
\label{sec:conclusion}

We presented a trace-driven framework for studying virtual PIM decision signals during LLM decoding and used it to characterize when offload value appears. The main conclusion is not that PIM-style offloading always helps, but that it becomes worthwhile only after the decode stream enters the right operating regime. In our current experiments, context length is the clearest coarse regime-entry variable, the onset is model-family dependent, and the observed boundary is consistent with rising KV-cache pressure. KV-aware features therefore look useful mainly as boundary-refinement signals on top of a context-centered characterization, especially near early-positive cases. At a higher level, our results suggest that the key scheduling problem in LLM decoding is not identifying kernel types, but detecting regime transitions driven by memory pressure.

The project has important limitations. Our backend is prior-work-grounded but still virtual rather than cycle-accurate, so we do not claim real hardware speedups. Our traces come from open-source runtimes rather than production serving engines, and runtime-stack diversity remains limited. Our model set, while broader than the initial study, still does not cover very large serving-scale models, extensive quantization variation, or family-within-size scaling. While the current workload set is limited in size, our goal is not to establish universal thresholds, but to demonstrate the existence and consistency of regime-dependent behavior across representative modern families. We therefore view the paper as an empirical characterization study rather than a final accelerator or scheduler paper.

The most useful next steps are to add a stronger runtime-diversity check, obtain at least one hardware or emulation anchor point for the virtual backend, and test whether the family-dependent onset patterns persist under broader serving conditions such as continuous batching.
